\section{Related Work}

\subsection{Many traditional depth estimation methods for light field cameras have been proposed, primarily exploiting physical cues such as Epipolar Plane Image (EPI) slopes, defocus, and photo-consistency correspondence [1]. These early approaches generally derive a theory based on the ideal Lambertian reflectance model, where the radiance of a 3D point remains constant across different viewpoints. Under this assumption, a spatial point projects as a straight line in the EPI, and the slope of this line is explicitly inversely proportional to the depth value [4]. While these methods perform well in extensive synthetic evaluation with ground-truth annotations, dealing with real-world non-Lambertian surfaces remains challenging.}

Building upon this geometric foundation, Wanner et al. [16] proposed a globally optimized depth estimation framework based on the structure tensor of EPIs. Specifically, they accurately calculated the local gradient orientations within the EPI to extract the line slopes, followed by a Markov Random Field (MRF) optimization to enforce spatial smoothness. This method successfully achieves high-quality depth maps in ideal environments with rich textures. However, it purely relies on the linear structure of the EPI, making it highly sensitive to noise and structural distortions caused by non-Lambertian reflections, and is mostly restricted to purely diffuse scenes.

While Wanner et al. focuses purely on EPI geometric structures, Tao et al. [17] extended the physical modeling by fusing defocus and correspondence cues. On the one hand, the defocus cue evaluates the sharpness of refocused images at different depth labels, which is significantly effective for texture-less regions. On the other hand, the correspondence cue measures the photo-consistency across angular views, performing well in textured areas. By adaptively combining these two complementary cues, their approach dramatically improves depth estimation in under-constrained regions. Nevertheless, when dealing with specular or translucent surfaces, both the defocus blur and the angular photo-consistency are severely corrupted, potentially leading to erroneous depth hypotheses and requiring pseudo labels for correction in many instances.

To further refine the correspondence matching and overcome the discrete angular sampling limitation, Jeon et al. [18] introduced a sub-pixel phase shift technique for light field depth estimation. By calculating the phase shift between sub-aperture images in the Fourier domain, they achieved sub-pixel accuracy in disparity estimation, up to sub-pixel precision, without requiring dense angular sampling. This simple yet effective strategy extensively improves the depth resolution and practical realization of light field cameras. However, the phase shift calculation inherently assumes constant brightness across views, which significantly fails when the intensity fluctuates dynamically due to view-dependent reflections, making its performance inferior to state-of-the-art methods in complex lighting conditions.

To address the violations of photo-consistency in complex scenes, Williem et al. [19] proposed a robust depth estimation method utilizing angular entropy and adaptive correspondence weights. They explicitly designed an angular entropy metric to detect occluded and non-Lambertian regions, subsequently down-weighting the unreliable correspondences in the cost aggregation step. Although this approach relatively improves robustness in mixed scenes by suppressing outliers, it attempts a binary or soft classification of reliable versus unreliable pixels. Such hand-crafted heuristics cannot physically model the complex light transport of non-Lambertian materials, often resulting in over-smoothed depth boundaries or missing details in highly reflective regions.

In short, despite their solid physical modeling in ideal scenarios, these traditional methods share critical limitations. First, they heavily rely on the Lambertian reflectance assumption and work poorly on non-Lambertian surfaces. When facing specular, translucent, or scattering materials, the EPI linear structures and defocus cues are substantially destroyed, causing the depth estimation to collapse dramatically. Second, the hand-crafted features and physical models exhibit poor generalization ability in complex mixed scenes. They cannot adaptively distinguish different material regions, and the iterative optimization processes incur high computational complexity, which is difficult to meet the requirements of practical applications. In contrast, our model explicitly addresses these issues through a novel unified dual-mask architecture. By proposing a three-layer angular signal decomposition mechanism, we reveal the fundamental reasons why non-Lambertian surfaces destroy the EPI linearity from both physical and geometric perspectives, providing a robust material identification scheme even under low angular resolution. Furthermore, our dual-mask routing mechanism effectively mitigates the depth collapse caused by the failure of physical assumptions, enabling high-accuracy depth estimation in both mixed and holistic scenes within a single unified model.

\subsection{Following the limitations of hand-crafted physical cues, deep learning-based methods have been extensively explored to implicitly learn geometric representations from light field data [20]. Early pioneering works primarily rely on Convolutional Neural Networks (CNNs) to extract local spatial and angular features. For instance, Shin et al. [10] proposed EPINet, which concatenates epipolar plane image (EPI)s (EPIs) and central view patches as inputs to a fully convolutional network, successfully regressing pixel-wise depth by learning the local slope variations. Building upon this geometric foundation, subsequent variants [21] introduced macropixel structures and multi-scale feature fusion to enhance the receptive field, significantly improving depth accuracy in textured regions. While these CNN-based architectures excel at capturing local geometric patterns, they often struggle with large displacements and textureless areas due to their inherently limited receptive fields.}

To address the limitations of local receptive fields in capturing global geometric consistency, recent approaches have incorporated attention mechanisms and Transformer architectures into light field depth estimation. While CNNs focus on local EPI slopes, Transformer-based models explicitly model long-range angular dependencies. Specifically, Liang et al. [22] proposed a Light Field Transformer that utilizes intra- and inter-view self-attention modules to aggregate global contextual information across multiple sub-aperture images. This design dramatically alleviates the ambiguity in textureless regions by enforcing global photo-consistency. Moreover, similar attention-driven frameworks [1] have been developed to dynamically weight reliable angular views, thereby suppressing noise from occluded boundaries. Despite their notable success in modeling global geometry, these methods generally treat all angular signals with equal physical importance, potentially amplifying errors when non-Lambertian reflections intrude.